% !PRINTED 734
\weberpage{734}

%<<CONT
and from this it follows that $S'_\beta$, just like $S_\beta$, has the same sign as all its conjugate values. But $S'_\beta$ arises from $S_\beta$ by putting $\mu$ in place of $m$, and by our assumption we therefore have

\[
e_0 \equiv 0,\ e_1 \equiv 0,\ \ldots,\ e_{\nu''-1} \equiv 0 \ (\mathrm{mod}\ 2),
\]

and, with the help of (19),

\[
e_{\nu''} \equiv 0,\ e_{\nu''+1} \equiv 0,\ \ldots,\ e_{\nu'-1} \equiv 0 \ (\mathrm{mod}\ 2),
\]

whereby the proof of Theorem 2. is carried through.

If we now choose in formula (12) the system $\tau_0, \tau_1, \ldots, \tau_{\nu'-1}$ for $\mathcal{E}_0, \mathcal{E}_1, \ldots, \mathcal{E}_{\nu'-1}$, and put the exponents $\xi$ into the form:

\[
\xi_0 = \frac{e_0}{e},\ \xi_1 = \frac{e_1}{e},\ \cdots,\ \xi_{\nu'-1} = \frac{e_{\nu'-1}}{e},
\]

in which $e, e_0, e_1, \ldots, e_{\nu'-1}$ are integral rational numbers without a common divisor, then we obtain

\[
\mathcal{E}(r)^e = \pm\, \tau_0^{e_0}\tau_1^{e_1} \cdots \tau_{\nu'-1}^{e_{\nu'-1}},
\]

and it follows from our theorem that $e$ must be \gesperrt{odd}; for if it were even, then $\mathcal{E}(r)^e$ would be the square of a real quantity, and consequently $\pm\, \mathcal{E}(r)^e$ would have the same sign as its conjugates. But then, by our theorem, all the $e_0, e_1, \ldots, e_{\nu'-1}$ would have to be divisible by 2, which again runs counter to the assumption that the exponents $e, e_0, \ldots, e_{\nu'-1}$ are to have no common divisor. We therefore have the theorem:

3. \gesperrt{If $\mathcal{E}(r)$ is any normal unit of the field $H_m$, then the odd number $e$ and the integers $e_0, e_1, \ldots, e_{\nu'-1}$ can be so determined that}

\weq{20}{\mathcal{E}(r_\beta)^e = \pm\, \tau_\beta^{e_0}\tau_{\beta+1}^{e_1} \cdots \tau_{\beta+\nu'-1}^{e_{\nu'-1}}}

\gesperrt{holds.}

This result leads to several important consequences:

4. \gesperrt{If $T$ is the regulator of the system of the $\tau_\beta$, and $T_0$ the regulator of a fundamental system of normal units in $H_m$, i.\,e. the minimal value which the regulator of any system of independent normal units can assume, then}

\weq{21}{T = C\,T_0, \quad C \equiv 1\ (\mathrm{mod}\ 2),}

\gesperrt{in which $C$ denotes an odd natural number.}
